% !PRINTED 434
\weberpage{434}
%<<CONT
und darin müssen, da $\tau$ vom 7\textsuperscript{ten} Grade sein soll, die $\varepsilon_1, \varepsilon_2, \varepsilon_3$
siebente Einheitswurzeln sein, die der Bedingung
\weq{2}{\varepsilon_1 \varepsilon_2 \varepsilon_3 = 1}
genügen.

Nun wollen wir die Substitution $\chi$ so zu bestimmen suchen,
dass die Bedingung $\chi \tau = \tau^4 \chi$ oder
\weq{3}{\tau \chi = \chi \tau^2}
erfüllt ist. Nehmen wir $\chi$ in der Form
\[
\chi = \begin{pmatrix} \alpha_1, & \alpha_2, & \alpha_3 \\ \beta_1, & \beta_2, & \beta_3 \\ \gamma_1, & \gamma_2, & \gamma_3 \end{pmatrix}
\]
an, so ergiebt die Bedingung (3)
\[
\begin{pmatrix} \alpha_1 \varepsilon_1, & \alpha_2 \varepsilon_1, & \alpha_3 \varepsilon_1 \\ \beta_1 \varepsilon_2, & \beta_2 \varepsilon_2, & \beta_3 \varepsilon_2 \\ \gamma_1 \varepsilon_3, & \gamma_2 \varepsilon_3, & \gamma_3 \varepsilon_3 \end{pmatrix}
=
\begin{pmatrix} \alpha_1 \varepsilon_1^2, & \alpha_2 \varepsilon_2^2, & \alpha_3 \varepsilon_3^2 \\ \beta_1 \varepsilon_1^2, & \beta_2 \varepsilon_2^2, & \beta_3 \varepsilon_3^2 \\ \gamma_1 \varepsilon_1^2, & \gamma_2 \varepsilon_2^2, & \gamma_3 \varepsilon_3^2 \end{pmatrix}.
\]

Da nun $\alpha_1, \beta_1, \gamma_1$ nicht alle drei verschwinden können, so
muss $\varepsilon_1^2$ mit einer der drei Grössen $\varepsilon_1, \varepsilon_2, \varepsilon_3$ übereinstimmen.
Wäre $\varepsilon_1 = \varepsilon_1^2$, also $\varepsilon_1 = 1$, und folglich $\varepsilon_3 = \varepsilon_2^{-1}$, so könnten
weder $\varepsilon_2$ noch $\varepsilon_3 = 1$ sein, weil sonst $\tau$ die identische Substitution
wäre. Es müsste also $\beta_1 = \gamma_1 = \alpha_2 = \alpha_3 = 0$ und $\varepsilon_2^2 = \varepsilon_3$,
$\varepsilon_2^3 = 1$ sein, was unmöglich ist. Es bleiben daher nach (2) noch
zwei mögliche Annahmen übrig:
\[
\varepsilon_1 = \varepsilon, \quad \varepsilon_2 = \varepsilon^2, \quad \varepsilon_3 = \varepsilon^4,
\]
\[
\varepsilon_1 = \varepsilon, \quad \varepsilon_2 = \varepsilon^4, \quad \varepsilon_3 = \varepsilon^2,
\]
worin $\varepsilon$ eine imaginäre siebente Einheitswurzel ist, und diese
beiden Annahmen führen zu zwei \gesperrt{verschiedenen} Gruppen, die
ganz gleichartig gebaut, übrigens auch isomorph sind. Wir verfolgen hier die erste Annahme weiter, setzen also
\weq{4}{\tau = \begin{pmatrix} \varepsilon, & 0, & 0 \\ 0, & \varepsilon^2, & 0 \\ 0, & 0, & \varepsilon^4 \end{pmatrix}, \quad \chi = \begin{pmatrix} 0, & 0, & 1 \\ 1, & 0, & 0 \\ 0, & 1, & 0 \end{pmatrix}, \quad \chi^2 = \begin{pmatrix} 0, & 1, & 0 \\ 0, & 0, & 1 \\ 1, & 0, & 0 \end{pmatrix}, \quad \chi^3 = 1.}

Dass wir die nicht verschwindenden Grössen $\alpha_3, \beta_1, \gamma_2$ gleich 1
gesetzt haben, ist keine Beschränkung, da wir jede andere Annahme durch eine Transformation der ganzen Gruppe darauf
zurückführen können.

Die Bedeutung von $\chi$ in dieser Form ist die einer cyklischen
Permutation der Variablen. Die Substitutionen $\tau, \chi$ erzeugen
zusammen eine Gruppe $\tau^\varrho \chi^\lambda$ vom 21\textsuperscript{sten} Grade.
